% Manuais do artefato. Incluir depois das referências bibliográficas.
\apendice[ape:manual_usuario]{Manual do usuário}

Este manual orienta o uso do protótipo Android do Luminthor. Os nomes entre aspas
correspondem aos controles da interface descrita no trabalho. As instruções não
substituem a validação de compatibilidade do APK distribuído em cada aparelho.

\section*{Idioma dos controles}

As mensagens de erro e as explicações são apresentadas em português, enquanto
as palavras reservadas da Lumi, como \texttt{if} e \texttt{print}, permanecem em
inglês. A interface utiliza ``Execute'' para executar, ``Trace'' para o depurador,
``Settings'' para as configurações e ``Templates'' para os modelos de código,
conforme a justificativa da Seção~\ref{sec:idioma_lumi}. As instruções abaixo
usam esses nomes para permitir localizar os controles.

\section*{Requisitos e instalação}

É necessário um aparelho Android compatível com o APK e espaço para o aplicativo
e os arquivos produzidos. O uso do APK não exige instalação de Flutter, Rust ou
ferramentas de desenvolvimento no celular. Uma conexão é necessária para obter o
instalador; a análise e a execução dos programas Lumi são realizadas pelo motor
local, sem envio do código a um servidor de compilação.

\begin{enumerate}
    \item Acesse o endereço de distribuição informado para a avaliação:
    \url{https://luminthor.vercel.app/}. Obtenha o arquivo APK disponibilizado pelo
    pesquisador. Se o endereço não estiver acessível, solicite o mesmo arquivo
    diretamente ao responsável pela distribuição.
    \item Abra o APK pelo navegador ou pelo gerenciador de arquivos do aparelho.
    Caso o Android solicite autorização para instalar aplicativos dessa fonte,
    siga a opção apresentada para o aplicativo que abriu o arquivo. Os nomes das
    opções variam conforme a versão do sistema e o fabricante.
    \item Confirme a instalação e abra o Luminthor. Aguarde a inicialização do
    armazenamento e do motor da linguagem; a animação de abertura pode ser pulada.
    \item Se a instalação ou a abertura falhar, registre o modelo do aparelho, a
    versão do Android, a identificação do APK recebido e a mensagem apresentada,
    sem incluir identificadores pessoais. Encaminhe esse registro ao pesquisador.
\end{enumerate}

O APK é específico para Android. Abrir o endereço de distribuição no navegador
não executa a IDE, e o arquivo APK não pode ser instalado em um iPhone. O requisito
mínimo do sistema é determinado pela compilação do APK; a configuração técnica
correspondente é descrita no Apêndice~\ref{ape:manual_tecnico}.

\section*{Criar, abrir, preservar e exportar arquivos}

Na tela inicial, selecione ``Novo arquivo'', informe um nome e confirme. A extensão
\texttt{.lumi} é acrescentada quando necessária. Os arquivos já utilizados aparecem
na seção ``Arquivos'' e nos recentes da Home. A busca permite localizar documentos
pelo nome. A seção ``Templates'' oferece modelos; ``Usar template'' abre uma cópia
para edição. ``Abrir arquivo'' importa um texto \texttt{.lumi} ou \texttt{.txt} pelo
seletor do aparelho. O conteúdo deve estar codificado em UTF-8; arquivos binários
ou em formato incompatível são rejeitados.

As alterações do editor são salvas automaticamente no armazenamento local do
aplicativo. Voltar à Home e reabrir o documento permite continuar a edição. Se a gravação falhar, o aplicativo exibe o aviso ``Não foi possível salvar o arquivo''; as alterações continuam abertas na sessão, e convém exportar o código para não perdê-lo. Esse salvamento é interno ao aplicativo. Para guardar uma cópia em outra pasta do celular, use o ícone ``Exportar arquivo'' na parte superior do editor, escolha o destino e confirme no seletor de arquivos.
O cancelamento do seletor não significa que uma cópia foi exportada.

O armazenamento é local, sem sincronização de documentos entre aparelhos. Para
preservar o trabalho antes de desinstalar o aplicativo ou limpar seus dados,
exporte os arquivos que deseja conservar.

\section*{Editar e executar um exemplo}

Abra um novo arquivo e insira o programa abaixo. A barra acima do teclado oferece
símbolos e palavras da linguagem; o autocompletar permite selecionar sugestões
durante a digitação. Use as configurações, acessíveis por ``Settings'', para
ajustar o tema e o tamanho da fonte. Quando o teclado ocultar a barra inferior,
feche-o para acessar as abas.

\begin{lstlisting}[language={},numbers=none]
var x = 2;
var y = 3;
x = x + y;
y = x * 2;
print(x);
print(y);
\end{lstlisting}

Selecione ``Execute''. A aba ``Saída'' do painel inferior deve apresentar 5 e 10,
nessa ordem e em linhas separadas. O estado final esperado é \texttt{x = 5} e
\texttt{y = 10}. Para experimentar outra entrada, altere o código e execute
novamente; a saída anterior não representa o resultado do programa modificado.

Os comandos utilizam palavras como \texttt{var}, \texttt{if}, \texttt{while},
\texttt{fn} e \texttt{print}. O comando \texttt{print} recebe uma expressão.
O ponto e vírgula é opcional, mas parênteses e chaves exigidos pelas construções
precisam estar completos. O protótipo não possui leitura interativa do teclado:
os valores de entrada dos exemplos devem ser alterados no próprio código.

\section*{Consultar diagnósticos, tokens e AST}

Se a execução indicar falha, abra ``Diagnósticos'' no painel inferior. Identifique
a categoria do erro, a linha e a coluna; leia o trecho destacado e a sugestão,
quando disponível. Corrija o código e selecione ``Execute'' outra vez. Uma falha
léxica, sintática ou semântica pode impedir a execução e, portanto, a produção
de saída ou de passos do depurador.

Para experimentar um erro sintático, execute o trecho abaixo. Falta o parêntese
de fechamento do comando de saída.

\begin{lstlisting}[language={},numbers=none]
var total = 10 + 5;
print(total;
\end{lstlisting}

Acrescente o parêntese, obtendo \texttt{print(total);}. A nova execução deve
produzir 15. Outro exemplo é declarar \texttt{var pontos = 8;} e tentar usar
\texttt{print(pntos);}: o nome não declarado gera um diagnóstico semântico,
que pode sugerir o identificador semelhante \texttt{pontos}.

A aba ``Tokens'' apresenta os elementos reconhecidos na análise léxica, suas
categorias e posições. A aba ``AST'' mostra a organização sintática do programa;
toque nos nós para expandir ou recolher seus filhos. A existência de tokens ou
de uma AST não significa, isoladamente, que o programa foi executado com sucesso.

\section*{Depuração passo a passo}

\begin{enumerate}
    \item Execute o primeiro exemplo deste manual, que produz um rastro válido.
    \item Abra ``Trace'', nome da aba do depurador, e selecione
    ``Iniciar Depuração Passo a Passo''.
    \item Use ``Avançar um passo'' e ``Voltar um passo'', ou a linha do tempo,
    para selecionar estados. Observe o número do passo, a linha, a ação e sua
    descrição.
    \item Consulte as variáveis e os tipos por escopo, a pilha de chamadas e a
    saída acumulada. No último passo, confira \texttt{x = 5}, \texttt{y = 10}
    e as duas linhas de saída.
    \item Selecione ``Parar depuração'' antes de editar o programa. Execute-o
    novamente para obter o rastro correspondente às alterações.
\end{enumerate}

Avançar ou retroceder consulta estados já registrados; não executa o programa
de novo. A edição permanece bloqueada durante a sessão de depuração. Quando não
há rastro, a aba informa que o código precisa ser executado ou que não foram
gerados passos. Um erro detectado antes da interpretação, como uma divisão cujo
divisor é comprovadamente zero, pode impedir qualquer rastro: \texttt{print(1 / 0);} e um divisor declarado com \texttt{const zero = 0;} são exemplos dessa verificação semântica. Já \texttt{var divisor = 0; print(1 / divisor);} falha durante a execução e pode preservar passos anteriores.

\section*{Trilha Lumi, limites e dificuldades de uso}

Em ``Aprender'', escolha a etapa disponível da Trilha Lumi. Cada aula apresenta
explicação, exemplo, desafio e pergunta de compreensão. ``Praticar no editor''
abre o arquivo associado à aula; retornar a ele preserva as edições anteriores.
Uma resposta correta à pergunta libera a próxima etapa. A prática no editor é
livre, sem correção automática do desafio, e repetir uma aula concluída não
duplica os pontos.

Se um programa continuar com erro, confira se todo o exemplo anterior foi
substituído, se os nomes coincidem e se os delimitadores estão completos. Uma
falha sintática deve ser corrigida antes de procurar o resultado da interpretação.
O motor limita a execução a 10.000 passos registrados; repetições prolongadas
devem ser revistas. Falhas de comunicação com a biblioteca nativa ou de abertura
do armazenamento exigem registrar a mensagem e comunicar o responsável pelo
aplicativo. Em problemas de importação ou exportação, confira o formato e o
destino escolhido antes de tentar novamente.

\apendice[ape:manual_tecnico]{Manual técnico de instalação e execução}

Este manual descreve a preparação do código-fonte e a geração do aplicativo.
Os comandos são executados em um terminal, a partir da raiz do repositório,
salvo indicação em contrário. As instruções refletem a configuração disponível
no projeto; sua reprodução deve registrar versões das ferramentas e plataforma
utilizada.

\section*{Ferramentas, dependências e plataformas}

O aplicativo utiliza Flutter/Dart e uma biblioteca nativa Rust. Embora a faixa
declarada em \texttt{pubspec.yaml} comece em Dart 3.3.0, o conjunto resolvido em
\texttt{pubspec.lock} exige Dart 3.13.0 ou superior, abaixo de 4.0.0, e Flutter
3.38.4 ou superior. O ambiente da verificação técnica documentada utilizou Flutter
3.47.1, Dart 3.13.1, rustc 1.96.0 e Cargo 1.96.0. 

As dependências de interface resolvidas incluem Riverpod 2.6.1, GoRouter 14.8.1,
Hive 2.2.3, Hive Flutter 1.1.0, \texttt{file\_picker} 13.1.0 e
\texttt{flutter\_svg} 2.3.0. O motor utiliza, entre outras, \texttt{logos}
0.16.1, \texttt{serde} 1.0.228 e \texttt{serde\_json} 1.0.149, conforme
\texttt{Cargo.lock}. Preserve os arquivos de resolução de dependências para
reproduzir a versão documentada.

O runtime Dart, a dependência Rust e o gerador de
\path{flutter_rust_bridge} devem usar a mesma versão exata: 2.12.0.
As ligações estão em \path{lib/src/rust} e
\path{rust_core/src/frb_generated.rs}; sua configuração está em
\path{flutter_rust_bridge.yaml}.

No Linux, são necessárias as ferramentas de compilação nativa e as bibliotecas
de desenvolvimento do GTK (\textit{GIMP Toolkit}, conjunto de ferramentas para interface gráfica). O README relaciona \texttt{clang}, \texttt{cmake},
\texttt{ninja-build}, \texttt{pkg-config}, \texttt{libgtk-3-dev},
\texttt{liblzma-dev}, \texttt{lld} e as ferramentas C/C++ da distribuição.
Os nomes dos pacotes e o gerenciador de instalação variam conforme o sistema.

\section*{Preparar e executar em Linux}

Depois de obter o código-fonte e instalar as ferramentas, verifique o ambiente
e resolva as dependências:

\begin{lstlisting}[language={},numbers=none]
flutter --version
rustc --version
cargo --version
flutter doctor
flutter pub get
cargo build --manifest-path rust_core/Cargo.toml
flutter devices
flutter run -d linux
\end{lstlisting}

A compilação Cargo produz \path{libide_engine.so} em
\path{rust_core/target/debug}. A configuração CMake do aplicativo Linux
também compila a biblioteca Rust e a inclui na pasta \texttt{lib} do pacote.
Para distribuir o aplicativo desktop, gere o pacote com
\texttt{flutter build linux} e preserve o diretório de distribuição completo,
incluindo bibliotecas e recursos. Copiar somente o executável não preserva
essas dependências.

Se houver conflito entre bibliotecas do ambiente de desenvolvimento e o sistema,
confira os erros de vinculação e a instalação de GTK e das ferramentas C/C++.
O README apresenta uma alternativa de execução com ambiente limpo para Linux.
Antes de alterar o código, registre a mensagem e confirme o ambiente usado.

\section*{Gerar o APK Android}

A configuração Android usa Gradle 8.14, Android Gradle Plugin 8.11.1, Kotlin
2.2.20 e alvo Java 17. O README orienta um kit de desenvolvimento Java (JDK, do inglês \textit{Java Development Kit}) 17 ou 21. Configure o kit de desenvolvimento de software Android (SDK, do inglês \textit{Software Development Kit})
em \path{android/local.properties}, pela chave \texttt{sdk.dir}, ou por
\path{ANDROID_HOME}/\path{ANDROID_SDK_ROOT}. O Flutter deve reconhecer
o SDK, o JDK e as ferramentas necessárias em \texttt{flutter doctor}.

Instale o kit de desenvolvimento nativo (NDK, do inglês \textit{Native Development Kit}) indicado pela propriedade \texttt{flutter.ndkVersion} do SDK
Flutter utilizado. O nível mínimo do aplicativo é herdado de
\texttt{flutter.minSdkVersion}, não fixado numericamente no repositório; no APK
publicado, corresponde à API 24 (Android 7.0), conforme a
Tabela~\ref{tab:medicoes_apk}. Não se deve confundir esse requisito com o nível
dos executáveis de compilação cruzada do Rust.

% A minipage descarta o espaço superior da listagem; o \medskip o repõe.
\medskip
\noindent\begin{minipage}{\linewidth}
\begin{lstlisting}[language={},numbers=none]
rustup target add aarch64-linux-android \
  armv7-linux-androideabi x86_64-linux-android
flutter build apk --release
\end{lstlisting}
\end{minipage}\par

A tarefa Gradle \texttt{buildRustAndroid} compila e empacota o motor para
\texttt{arm64-v8a}, \texttt{armeabi-v7a} e \path{x86_64}. Sua configuração
atual procura a ferramenta NDK de um hospedeiro Linux x86\_64. A geração em outro sistema hospedeiro exige adaptar essa etapa.
O APK fica em \path{build/app/outputs/flutter-apk/app-release.apk}.

O perfil \texttt{release} usa a chave de assinatura de depuração na configuração
atual. Isso identifica um artefato de protótipo e não uma configuração de
publicação em loja. Antes de distribuir uma nova compilação, registre a revisão
do código, as ferramentas, a versão, o arquivo gerado e seu hash. A atualização
de uma instalação existente precisa manter uma assinatura compatível.

\section*{Regenerar a ponte e verificar o motor}

Regenerar as ligações é necessário quando a API Rust exposta ao Flutter
for alterada. Para isso, instale o gerador correspondente e execute, na raiz:

\begin{lstlisting}[language={},numbers=none]
cargo install flutter_rust_bridge_codegen --version 2.12.0
flutter_rust_bridge_codegen generate
cargo build --manifest-path rust_core/Cargo.toml
\end{lstlisting}

O motor pode ser experimentado sem abrir a interface, pela ferramenta
\texttt{ide-cli}. No exemplo abaixo, a saída do programa deve ser 5; a opção
\texttt{--json} inclui também tokens, AST, diagnósticos e rastro disponíveis.

\begin{lstlisting}[language={},numbers=none]
cargo run --manifest-path rust_core/Cargo.toml --bin ide-cli \
  -- --json "var x = 2; x = x + 3; print(x);"
\end{lstlisting}

\section*{Testes, diagnóstico técnico e limites de reprodução}

Para a verificação automatizada disponível no repositório, execute:

\begin{lstlisting}[language={},numbers=none]
cargo test --manifest-path rust_core/Cargo.toml
flutter analyze
flutter test
\end{lstlisting}

Os testes de FFI no Flutter carregam a biblioteca nativa compilada; por isso,
\texttt{cargo build} deve antecedê-los. Registre o resultado dos comandos na
mesma revisão de código. Uma falha ao carregar o motor deve levar à conferência
da biblioteca, da arquitetura e da versão da ponte. Erros do compilador Lumi
devem ser examinados pela categoria e pela posição retornadas, preservando o
programa que os reproduz.

O procedimento documentado cobre a execução em Linux e a geração do APK Android; as demais pastas de plataforma do projeto Flutter não foram preparadas. A execução no navegador exigiria uma adaptação do motor para WebAssembly. 

\section*{Registro da verificação final}

A versão final descrita neste trabalho corresponde ao commit
\hashcod{33a64fc3b792c3ee26d7a69ea0c6fb3205aeacd3}. Os marcos do
Quadro~\ref{quad:iteracoes} correspondem aos commits
\hashcod{0d0f7191fc07aaad7cce24cbbc51aea8c8003cac} (base do pacote de 10/10/2026,
acrescido de alterações locais registradas no manifesto do checkpoint) e
\hashcod{d2fbb0258ea2b4b21d95baac6796a1d8a83590ac} (11/11/2026). A detecção semântica
de divisão por zero foi incorporada no commit
\hashcod{84397f8620f797de06264c58c275bb2103903667} (12/11/2026).

A verificação foi executada em Fedora Linux x86\_64, com Flutter 3.47.1, Dart
3.13.1, \texttt{rustc} e Cargo 1.96.0. A cobertura de ramos do motor exige o
compilador \textit{nightly} (\texttt{rustc} 1.102.0-nightly) e a ferramenta
\texttt{cargo-llvm-cov} 0.9.1. Os comandos utilizados, a partir da raiz do
repositório, foram:

\begin{lstlisting}[language={},numbers=none]
cargo build --offline --manifest-path rust_core/Cargo.toml
cargo test --offline --manifest-path rust_core/Cargo.toml
flutter analyze --no-pub
flutter test --no-pub --coverage
cargo +nightly llvm-cov --manifest-path rust_core/Cargo.toml \
  --branch --lib --tests \
  --ignore-filename-regex 'frb_generated|bin/cli'
cargo run --release --example medicoes \
  --manifest-path rust_core/Cargo.toml
cargo run --example reproducao_t3 \
  --manifest-path rust_core/Cargo.toml
\end{lstlisting}

A cobertura Dart foi calculada a partir de \path{coverage/lcov.info}; o recorte sem
ligações geradas exclui a pasta \path{lib/src/rust/}. Na cobertura Rust, foram
excluídos o arquivo de ligações gerado (\path{frb_generated.rs}) e a CLI
(\path{src/bin/cli.rs}), que não possui testes. Os programas
\path{rust_core/examples/medicoes.rs} e \path{rust_core/examples/reproducao_t3.rs}
reproduzem, respectivamente, as medições da Tabela~\ref{tab:medicoes_motor} e os
cenários do Quadro~\ref{quad:reproducao_t3}.

O APK da versão final foi gerado com \texttt{flutter build apk --release} e lido com
\texttt{aapt dump badging} (Android \textit{build-tools} 36.1.0). Seu resumo
criptográfico SHA-256 (do inglês \textit{Secure Hash Algorithm}, com 256 bits) é
\hashcod{c1c7cd65d1766a3c3ece65af204191b8d68fffba59afec2c7cceb14c0aacb19d}.
